%%% Evaluation.tex --- 
%% 
%% Filename: Evaluation.tex
%% Author: Jacob Ringfjord & Max Wilén
%%% Code:

\chapter{Evaluation}
\label{cha:evaluation}

Assuming that some kind of prototype/artifact/system/component has been developed/adapted/changed/implemented this should typically be evaluated. This chapter then becomes a ``mini''-version of the IMRaD model. The structure below is just one example of how one can structure this kind of chapter. It must be adapted to the kind of evaluation that is being done.

\section{Evaluation method}

Describe the method used in the evaluation. Be detailed and specific. Use tables to summarise parameters used if you are running simulations. Describe what experiments that are conducted, and their purpose. Explain the metrics used, and motivate their choice.

\section{Results}

Show the results of your evaluation. Explain to the reader what the results show.

\section{Summary of results}

It is often useful to provide some kind of summary of the results, especially if there are many graphs and tables. 

%%%%%%%%%%%%%%%%%%%%%%%%%%%%%%%%%%%%%%%%%%%%%%%%%%%%%%%%%%%%%%%%%%%%%%
%%% Evaluation.tex ends here
